La corrida definitiva contiene \textbf{1800} unidades experimentales y
30 repeticiones por celda. Para $N=100000$, la
correlación de rangos entre $\gamma$ y migraciones por inserción, promediando
$\tau$, fue -1.000; entre $\tau$ y colisiones por pares por clave,
promediando $\gamma$, fue 0.949. El cociente medio entre pares
observados y la cota universal fue 0.964. Estos valores describen
la muestra de semillas y no convierten una esperanza en cota determinista.

La menor migración media apareció en $(\gamma,\tau)=(4,
0.70)$; el menor uso estructural final en
(1.5,0.90); y la menor tasa de pares en
(4,0.50). Que estos puntos no
sean necesariamente el mismo tratamiento aporta evidencia empírica del
compromiso multiobjetivo.

\begin{table}[ht]
\centering\small
\caption{Promedios por $\gamma$ para $N=100000$, agregados sobre $\tau$ y semillas.}
\begin{tabular}{rrrrrr}
\toprule
$\gamma$ & migr./ins. & bytes/elem. & pares/clave & resizes & ns/oper. \\
\midrule
1.25 & 4.508 & 157.7 & 0.333 & 40.2 & 66080.8 \\
1.5 & 2.413 & 166.1 & 0.289 & 23.0 & 43657.2 \\
2 & 1.491 & 187.3 & 0.222 & 14.0 & 33174.9 \\
3 & 0.929 & 218.5 & 0.204 & 9.0 & 29056.5 \\
4 & 0.723 & 276.8 & 0.189 & 7.2 & 27944.1 \\
\bottomrule
\end{tabular}
\end{table}
